\chapter{Einleitung}
\label{ch:Einleitung}

%
% Was ist ein APT-Scanner?
%
Ein APT-Scanner (oder präziser ein APT-Angriffserkennungssystem)
ist eine spezialisierte Sicherheitssoftware,
die darauf ausgelegt ist,
die Aktivitäten von fortgeschrittenen andauernden Bedrohungen zu identifizieren.
Während herkömmliche Virenscanner oft auf bekannten Signaturen basieren,
müssen APT-Scanner deutlich komplexere Methoden anwenden,
um die typischerweise „subtilen"
und über lange Zeiträume gestreckten Angriffe zu finden.
Der Begriff APT steht für Advanced Persistent Threat
und bezeichnet eine Form des Cyberangriffs,
bei der sich ein Angreifer unbefugten Zugang zu einem Computersystem verschafft
und potenziell über einen langen Zeitraum unentdeckt darin verbleibt.
Im Gegensatz zu herkömmlichen Angriffen
handelt es sich bei APTs um menschliche Akteure
mit hoher Motivation und hohen Ressourcen,
nicht um einfachen, automatisierten Schadcode.

%
% Anomalieerkennung
%
Da APT-Akteure häufig neue Schwachstellen (Zero-Day-Exploits) nutzen
oder legitime Systemwerkzeuge missbrauchen (Living-off-the-Land),
helfen einfache Blacklists oft nicht weiter.
Moderne APT-Angriffserkennungssysteme nutzen daher Anomalieerkennung.
Die Anomalieerkennung identifiziert ungewöhnliche Muster
oder Instanzen in Daten,
die nicht dem erwarteten Verhalten entsprechen.
Sie unterscheidet im Gegensatz zu herkömmlichen Virenscannern
nicht zwischen bekannten und unbekannten Angriffen.
Da die Anomalieerkennung auf dem Systemverhalten basiert,
um Abweichungen zu erkennen,
sind APT-Angriffserkennungssysteme
mit einem deutlich höheren Wartungsaufwand verbunden.
Das APT-Angriffserkennungssystem lernt zunächst
das „normale" Verhalten des Rechners.
Jede signifikante Abweichung wird als potenzieller Angriff gemeldet.
Analysten bewerten die potenziellen Angriffe
und erkennen False Positives,
welche ab diesem Zeitpunkt automatisch
von dem APT-Angriffserkennungssystem als solche eingeordnet werden.

%
% Performance-Relevanz und Ziel der Arbeit
%
APT-Angriffserkennungssysteme werden auf einer hohen Anzahl
von produktiven Systemen,
häufig auf kritischen Servern
oder in ressourcenbeschränkten Umgebungen, eingesetzt.
Hierbei werden regelmäßig Scans ausgeführt.
Auch bei Incident-Response-Szenarien
kommen APT-Angriffserkennungssysteme zum Einsatz.
Die Performance des APT-Angriffserkennungssystems
ist von entscheidender Bedeutung.
Lange Scan-Zeiten und hohe Ressourcenauslastung
stellen ein relevantes Problem dar.
Das Risiko von Performanceeinbußen während des Scans
auf produktiven Systemen muss möglichst gering gehalten werden.
Gleichzeitig muss ein Scan auf einem potenziell befallenen System
eine niedrige Scan-Zeit aufweisen
und die zur Verfügung stehenden Ressourcen
möglichst weitreichend nutzen.
Um diese Voraussetzungen an die Performance zu erfüllen,
ist eine effiziente Parallelisierung durch einen Scheduler
von hoher Relevanz.
In dieser Arbeit wird der Versuch vorgenommen,
die Performance des APT-Scanners Thor
der Firma Nextron Systems GmbH,
im Folgenden als Scanner bezeichnet, zu optimieren,
indem ein task-basierter Scheduler implementiert wird.

\section{Grundlagen}
In diesem Kapitel wird auf die Arbeitsweise des Scanners eingegangen, insofern es für das Verständnis dieser Arbeit notwendig ist.
Es wird kurz auf den aktuellen Stand der Forschung im Bezug auf Scheduling Algorithmen eingegangen.

Was ist Performance?




Die drei Bestandteile des Namens beschreiben die Art der Bedrohung genau:

Bei Advanced-Persistent-Thread-Scannern (APT-Scanner), die unter anderen in Incidence-Response-Szenarien eingesetzt werden, ist die Performance sehr relevant.
Aufgrund des Einsatzes auf einer hohen Anzahl produktiver Systeme, häufig auf kritischen Servern oder ressourcenbeschränkten Umgebungen stellen
lange Scan-Zeiten und hohe Ressourcenauslastung ein relevantes Problem dar.
Sowohl die Laufzeit von Scans bei Incident-Response-Maßnahmen
als auch das Risiko von Performance Einbußen während Scans auf produktiven Systemen sollten möglichst gering gehalten werden.
Vor diesem Hintergrund ist eine Optimimierung der Scanner-Architektur hinsichtlich Parallelisierung, Ressourcennutzung und Durchsatz sowohl aus ökonomischer
als auch aus sicherheitstechnischer Sicht von hoher Relevanz.

In dieser Arbeit wird anhand des APT-Scanners Thor der Firma Nextron Systems, im folgenden als Scanner bezeichnet, der Frage nachgegangen
ob die Performance durch die Implementierung von task-basierten Scheduling verbessert werden kann.
Die Hypothese lautet, dass die Implementierung des task-basierten Schedulers die Performance des Scanners im Bezug auf die Hauptspeicherauslastung oder
die Scan-Zeit bei der Verwendung von mehreren Threads verbessern kann.

\section{Grundlagen}
In dieser Sektion wird für das Verständnis der Arbeit nötiges Wissen zu diesen Themen vermittelt:
\begin{enumerate}
  \item Bisherige Arbeitsweise des Scanners
  \item Lösungsansatz
  \item Stand der Forschung
  \item Überblick über die Arbeit
\end{enumerate}

\subsection{Bisherige Arbeitsweise des Scanners}



\section{}
%
% Section: Der erste Abschnitt
%
\section{Der erste Abschnitt des Kapitels}
\label{sec:background:first_section}
Sia ma sine svedese americas. Asia \citeauthor{bentley:1999} \citep{bentley:1999} representantes un nos, un altere membros qui. De web nostre historia angloromanic. Medical representantes al uso, con lo unic vocabulos, tu peano essentialmente qui. Lo malo laborava anteriormente uso.

\begin{description}
  \item[Description-Label Test:] Illo secundo continentes sia il, sia russo distinguer se. Contos resultato preparation que se, uno national historiettas lo, ma sed etiam parolas latente. Ma unic quales sia. Pan in patre altere summario, le pro latino resultato.
  \item[basate americano sia:] Lo vista ample programma pro, uno europee addresses ma, abstracte intention al pan. Nos duce infra publicava le. Es que historia encyclopedia, sed terra celos avantiate in. Su pro effortio appellate, o.
  \item[Cras venenatis:] Purus et posuere lacinia, nisl sapien dapibus metus, a ornare enim odio in ipsum. Quisque imperdiet nibh metus, in fringilla tellus. Duis varius dui eget orci commodo ac sollicitudin est placerat. Cras varius tincidunt arcu, quis imperdiet nibh rhoncus vel. Sed non justo orci, non accumsan felis. Maecenas condimentum convallis.
\end{description}
Tu uno veni americano sanctificate. Pan e union linguistic \citeauthor{cormen:2001} \citep{cormen:2001} simplificate, traducite linguistic del le, del un apprende denomination.

\subsection{Ein Unterabschnitt}
\label{subsec:background:first_section:first_subsection}
Uno pote summario methodicamente al, uso debe nomina hereditage ma. Iala rapide ha del, ma nos esser parlar. Maximo dictionario sed al. Aenean posuere, enim in ultricies facilisis, ligula lacus eleifend eros, accumsan commodo metus justo placerat justo. Donec sit amet mauris dolor, at imperdiet lacus. In laoreet pretium condimentum. Proin ut varius diam. Fusce ipsum ipsum, elementum id porttitor at, pharetra congue nisi.

\subsection{Ein weiterer Unterabschnitt}
\label{subsec:background:first_section:second_subsection}
Deler utilitate methodicamente con se. Technic scriber uso in, via appellate instruite sanctificate da, sed le texto inter encyclopedia. Ha iste americas que, qui ma tempore capital. Class aptent taciti sociosqu ad litora torquent per conubia nostra, per inceptos himenaeos. Proin vitae urna id metus vestibulum lobortis. Duis rhoncus pulvinar massa, eget venenatis justo dapibus sed.

%
% Section: Der Zweite Abschnitt
%
\section{Ein zweiter Abschnitt}
\label{sec:background:second_section}
Phasellus ut ipsum nulla, vitae venenatis augue. Suspendisse potenti. Mauris suscipit justo a dolor laoreet lacinia. Pellentesque habitant morbi tristique senectus et netus et malesuada fames ac turpis egestas. Aliquam commodo commodo dui, nec auctor mi malesuada et. Aenean tortor erat, semper eu ullamcorper non, dignissim sed lectus. Praesent et pretium leo.

\subsection{Ein Unterabschnitt}
\label{subsec:background:second_section:first_subsection}
Vivamus at massa ut turpis dignissim mattis. Vivamus odio metus, venenatis vitae malesuada et, dignissim sed nunc. Mauris a nisl id massa viverra mattis in ultrices odio. Vestibulum ante ipsum primis in faucibus orci luctus et ultrices posuere cubilia Curae; Curabitur quis metus ac sem venenatis dignissim nec.

\subsubsection{Ein Unter-Unterabschnitt}
\label{ssubsec:background:second_section:first_subsection:first_subsubsection}
Sed vel ante vel quam commodo cursus. Class aptent taciti sociosqu ad litora torquent per conubia nostra, per inceptos himenaeos. Duis non turpis eget quam rutrum scelerisque. Duis nec quam metus. Curabitur purus dui, sagittis vel mattis a, elementum vitae risus. Pellentesque a tellus lacus, id gravida lectus.

